% Chapter 2, Figure 19: overdamped response to a step disturbance in yaw (scan: figures/ch2/fig19.png).
% Curve: fig19.py, eq. (35) with eqs. (25), (36a), (36b), zeta = 2.1 (matched to the scan on the family time axis),
% in units of M_s/C_1 against omega_n t; the dashed line is the asymptote M_s/C_1, approached from below.
% One layout and time scale for Figs 16-19 (same axes, t to 2.75 pi/omega_n): Fig 18 has the same C_1/I_L,
% so this curve starts with the same curvature and then rises more slowly (the caption).
\documentclass[11pt]{standalone}
\usepackage{tamrfig}
\begin{document}
\begin{tikzpicture}
  \begin{axis}[tamr sketch,
      width=3.8in, height=2.5in,
      xmin=0, xmax=9.35, ymin=-0.36, ymax=2.62,
      ytick={1}, yticklabels={$\dfrac{M_s}{C_1}$},
      xlabel={$t$ (sec)}, ylabel={$\alpha_X$ (rad)},
      % tamr sketch's rotate=-90 turns this label to read downward: reset it to upright, above the arrow
      every axis y label/.style={at={(ticklabel* cs:1.0)}, anchor=south}]
    \draw[guide] (axis cs:0,1) -- (axis cs:8.63938,1);
    \addplot[series1] table[x=t, y=alpha, col sep=comma] {figures/v2/ch2/fig19.csv};
    \node[anchor=east, font=\footnotesize, text=ink2] at (axis cs:0,0) {0};
  \end{axis}
\end{tikzpicture}
\end{document}
